\section{Introduction}
\label{sec:intro}

The Hubble tension is an observational disagreement between the local
supernova distance ladder,
$H_{0} = 73.04 \pm 1.04$\kmsMpc{} \citep{Riess2022}, and the Planck
2018 TT/TE/EE spectrum, $67.4 \pm 0.5$\kmsMpc{} in flat $\Lambda$CDM
\citep{Planck2018}, reviewed in \citet{Verde2019} and
\citet{KnoxMillea2020}. The mismatch extends beyond a single number:
DESI DR2 BAO \citep{DESIDR2} sharpens the late-time
distance--redshift relation while the Pantheon+SH0ES compilation
\citep{Scolnic2022,Brout2022} maps the supernova expansion history,
and neither a distance-calibrated nor a CMB-calibrated history fits
all sectors at once. Observed values, uncertainties, and covariance
matrices are held fixed across model comparisons; only theory
predictions change between branches.

The standard question this paper asks is: is the disagreement
reducible to a compact redshift-dependent difference between
independently calibrated background histories? Rather than fitting
another cosmological model with additional dynamical freedom, we test
whether the observable gap between already-calibrated endpoints has
low-dimensional, predictable structure.

The experimental design uses three frozen endpoints. The
distance-compatible history~R reaches
$H_{0}=\DistanceHZero$\kmsMpc{} and fits DESI DR2 BAO and
Pantheon+SH0ES supernovae. The CMB-compatible history~C1 is a Planck TT/TE/EE-compatible
reference with $H_{0}=\CmbOneHZero$\kmsMpc{} produced with the stated
CAMB/likelihood configuration. The secondary Planck-compatible reference~C2 is a
separately implemented baseline with
$H_{0}=\CmbTwoHZero$\kmsMpc{}, built through its own likelihood and
spectral-parity path for reference-sensitivity checking. C1 and C2 are separately implemented in pipeline/model space, not statistically
independent in data space: they share the
Planck sky but differ in construction pipeline, parameter choices,
and calibration route.

This differs from fitting another cosmological model in four ways.
First, the endpoints are constructed separately, before any
comparison is made. Second, bridge dimensionality is selected with
held-out prediction, not training $\chi^{2}$. Third, higher
flexibility is allowed to fail: overfitting orders and
boundary-limited families are retained as frozen adversarial
controls. Fourth, the reference-branch robustness test changes the CMB
reference branch without changing the hypothesis: the same compact
map must describe both $\Delta_{1}$ and $\Delta_{2}$.

The headline result, over the constrained domain $z\le1.8$, is
\[
\Delta_{1}(z) \simeq \Delta_{2}(z),
\qquad
\|\Delta_{C}\| \ll \|\Delta_{1,2}\|,
\]
where the fields
$\Delta_{1}(z)\equiv\ln H_{R}(z)-\ln H_{C_{1}}(z)$,
$\Delta_{2}(z)\equiv\ln H_{R}(z)-\ln H_{C_{2}}(z)$, and
$\Delta_{C}(z)\equiv\ln H_{C_{1}}(z)-\ln H_{C_{2}}(z)$
are defined in Section~\ref{sec:endpoint-repr} and summarized in the
nomenclature table (Table~\ref{tab:nomenclature}).
The two R--CMB difference fields share one dominant direction
(cosine~\GeometryCosineSimilarity{}, principal angle
\GeometryPrincipalAngle$^{\circ}$), peak near
$z\approx\GeometryPeakZ{}$, and transfer across branches with
fractional residual degradation~\GeometryTransferDegradation{},
whereas the CMB--CMB difference is much smaller
throughout (about 1--8\% of the R--CMB difference across the
constrained bins).

In plain terms, the distance-sector endpoint requires a
systematically higher expansion rate than the Planck-calibrated
references, and the size of that mismatch changes with redshift
rather than behaving like a simple constant $H_{0}$ offset. A compact
five-parameter map summarizes that pattern over $z\le1.8$, but its
statistical significance and physical origin are not established, and
the same compact extrapolation fails at the Ly$\alpha$ point at
$z=2.33$.

This paper identifies and characterizes that empirical deformation.
The principal result of this work is therefore not a specific
dynamical model, but an empirically defined expansion-history target
that candidate explanations of the CMB--distance discrepancy can be
tested against. Its physical origin remains deliberately unaddressed.
Section~\ref{sec:endpoints} defines the data and the R, C1, and C2
endpoints (with the nomenclature table immediately below);
Section~\ref{sec:methods} the empirical reconstruction method;
Section~\ref{sec:bridge} the predictive model-selection tests;
Section~\ref{sec:robustness} the reference-branch robustness test;
Section~\ref{sec:results} the results;
and Sections~\ref{sec:discussion}--\ref{sec:conclusion} the
interpretation, limitations, and conclusions.
